\chapter{Mesaj Protokolleri}
\label{app:protokol}

Bu ek, sistemde tanımlanan tüm UART ve BLE mesaj çerçevelerini tek yerde
listeler. Mesajlar üç hat üzerinde dolaşmaktadır:

\begin{itemize}[itemsep=1pt]
    \item \textbf{UART1} (115200~bps, 8N1): STM32 $\leftrightarrow$ ESP32-CAM
    \item \textbf{UART2} (9600~bps, 8N1): STM32 $\leftrightarrow$ HM-10
    \item \textbf{BLE GATT} (FFE0/FFE1): HM-10 $\leftrightarrow$ Telefon
          (transparent, UART2'nin baytları aynen iletilir)
\end{itemize}

\section{UART1 --- STM32 / ESP32-CAM}

\paragraph{STM32 $\rightarrow$ ESP32-CAM}

\begin{tabular}{l l}
\toprule
\textbf{Komut} & \textbf{Anlam}\\
\midrule
\texttt{CAPTURE\textbackslash n}    & 10-frame burst başlat ve sunucuya gönder\\
\bottomrule
\end{tabular}

\paragraph{ESP32-CAM $\rightarrow$ STM32 (RESULT)}

Burst tamamlandığında veya hata oluştuğunda tek satırlık CSV:

\begin{tabular}{l l}
\toprule
\textbf{Satır}                              & \textbf{Anlam}\\
\midrule
\texttt{1;<name>;<sim>\textbackslash n}     & MATCH; isim ve similarity\\
\texttt{0;;<sim>\textbackslash n}           & NOMATCH; en yakın similarity\\
\texttt{ERR:<tag>\textbackslash n}          & Hata; tag örnek: no\_wifi, http\_500, json\_parse, timeout\\
\bottomrule
\end{tabular}

\section{UART2 ve BLE --- STM32 / HM-10 / Telefon}

HM-10 transparent çalıştığı için STM32'nin UART2'ye yazdığı satırlar BLE
notification olarak telefona iletilir.

\begin{tabular}{l l}
\toprule
\textbf{Satır}                                   & \textbf{Anlam}\\
\midrule
\texttt{SCANNING\textbackslash n}                & TARA başlangıcı\\
\texttt{MATCH:<name>;<sim>\textbackslash n}      & Pozitif eşleşme (telefon 155'i arar)\\
\texttt{NOMATCH\textbackslash n}                 & Negatif burst\\
\texttt{PANIC\textbackslash n}                   & PANİK butonu (telefon 155'i arar)\\
\texttt{NETERR\textbackslash n}                  & Scan timeout veya transport hata\\
\texttt{HB\textbackslash n}                      & 5\,s aralıkla kalp atışı (yalnızca IDLE'da)\\
\bottomrule
\end{tabular}

\section{HTTP --- ESP32-CAM / Sunucu}

\paragraph{İstek.}

\begin{lstlisting}[language=none, basicstyle=\ttfamily\footnotesize]
POST /search HTTP/1.1
Host: 18.192.45.175:8000
Content-Type: application/octet-stream
X-Session-Id: 7c5f9a3b-...-e1
X-Frame-Index: 7
X-Frame-Total: 10
X-Shared-Secret: <token>     # opsiyonel

<JPEG bytes>
\end{lstlisting}

\paragraph{Ara kare yanıtı.}

\begin{lstlisting}[language=none, basicstyle=\ttfamily\footnotesize]
HTTP/1.1 200 OK
Content-Type: application/json

{"status":"continue","frames_received":7,"frames_total":10,"quality_ok_this_frame":true}
\end{lstlisting}

\paragraph{Son kare yanıtı.}

\begin{lstlisting}[language=none, basicstyle=\ttfamily\footnotesize]
HTTP/1.1 200 OK
Content-Type: application/json

{
  "match": true,
  "name": "Ali_Yilmaz",
  "similarity": 0.9437,
  "frames_used": 10,
  "frames_total": 10,
  "liveness_score": 0.014823
}
\end{lstlisting}

\paragraph{Hata yanıtları.}

\begin{tabular}{l l l}
\toprule
\textbf{Kod} & \textbf{Body} & \textbf{Anlam}\\
\midrule
400 & \texttt{\{"error":"bad\_index\_headers"\}}    & X-Frame-* başlıkları geçersiz\\
400 & \texttt{\{"error":"missing\_session\_id"\}}   & X-Session-Id yok veya boş\\
400 & \texttt{\{"error":"bad\_frame\_indices"\}}    & idx > total veya $\leq$ 0\\
400 & \texttt{\{"error":"image\_too\_small"\}}      & JPEG $<$ 1\,KB\\
401 & \texttt{\{"error":"unauthorized"\}}           & X-Shared-Secret eşleşmedi (auth aktif)\\
500 & \texttt{\{"error":"session\_lost"\}}          & Oturum timeout sonrası temizlenmiş\\
\bottomrule
\end{tabular}
